\subsection{Zero-curvature equation and the Landau--Lifshitz equations of motion}
\label{sec:class6-eom}

We derive the continuum Hamiltonian flow from the Class 6 bond Hamiltonian
and compare it with the corrected Lax pair of
Section~\ref{sec:class6-time-lax}.  The Hamiltonian also fixes the
continuum limit of the Heisenberg derivative on the left-hand side of
the lattice zero-curvature equation.

With $\boldsymbol S_n(t)=\boldsymbol S(x_n,t)$ and
$\kappa_6=\epsilon^2\alpha_6$, expanding the bond lower symbol in
Eq.~\eqref{eq:class6-bond-lower-symbol} gives
\begin{equation}
h_{6;n,n+1}^\downarrow
=2+\epsilon^2\mathcal H_6(x_n,t)+\mathcal O(\epsilon^3),
\label{eq:class6-h-lower-expansion}
\end{equation}
where the continuum Hamiltonian density at $(x,t)$ is
\begin{equation}
\mathcal H_6(x,t)
=-\frac12\,\partial_x\boldsymbol S^{\mathsf T}\partial_x\boldsymbol S
+\alpha_6 S^+S^3.
\label{eq:class6-hamiltonian-density}
\end{equation}
Here we used the spatial derivatives of the unit-spin constraint in
Eq.~\eqref{eq:spin-kinematical-identities}.  The anisotropy can be written
in terms of the constant symmetric matrix
\begin{equation}
N_6:=\begin{pmatrix}
0&0&-1\\
0&0&-\ii\\
-1&-\ii&0
\end{pmatrix},
\qquad N_6^{\mathsf T}=N_6,
\qquad N_6^3=0.
\label{eq:class6-N-definition}
\end{equation}
Since $\boldsymbol S^{\mathsf T}N_6\boldsymbol S=-2S^+S^3$, the
continuum Hamiltonian at fixed circumference $\ell=N\epsilon$ is
\begin{align}
H^{\rm cl}_6
&:=\lim_{\epsilon\to0}
\frac{\bigl(H^{\rm lat}_6\bigr)^\downarrow-2N}{\epsilon}
\nonumber\\
&=\int_0^\ell dx\left[
-\frac12\,\partial_x\boldsymbol S^{\mathsf T}\partial_x\boldsymbol S
-\frac{\alpha_6}{2}\boldsymbol S^{\mathsf T}N_6\boldsymbol S
\right].
\label{eq:class6-continuum-hamiltonian}
\end{align}

We use the equal-time spin Poisson bracket with the same normalization as in
Eq.~\eqref{eq:class5-poisson-bracket},
\begin{equation}
\{S^i(x),S^j(y)\}_6
=2\,\varepsilon_{ijk}S^k(x)\,\delta(x-y),
\qquad \varepsilon_{123}=1,
\label{eq:class6-poisson-bracket}
\end{equation}
where $\varepsilon_{ijk}$ is completely antisymmetric and $\delta(x-y)$
is the Dirac delta function.  The common time argument is suppressed
in Eq.~\eqref{eq:class6-poisson-bracket}.
We take the functional derivative with respect to the Cartesian spin
components before imposing $\boldsymbol S^{\mathsf T}\boldsymbol S=1$.
Integration by parts with periodic variations then gives
$\delta H^{\rm cl}_6/\delta\boldsymbol S
=\partial_x^2\boldsymbol S-\alpha_6N_6\boldsymbol S$.
The Hamiltonian equation
$\partial_tS^i=\{S^i,H^{\rm cl}_6\}_6$ therefore reads
\begin{equation}
\partial_t\boldsymbol S
=-2\,\boldsymbol S\times
\left(\partial_x^2\boldsymbol S-\alpha_6N_6\boldsymbol S\right).
\label{eq:class6-vector-eom}
\end{equation}
The cross product is bilinear, as in
Section~\ref{sec:class5-eom}, and the right-hand side preserves
$\boldsymbol S^{\mathsf T}\boldsymbol S=1$.
In terms of $S^\pm$ and $S^3$, Eq.~\eqref{eq:class6-vector-eom} becomes
\begin{align}
\partial_tS^+
&=2\ii\left(
S^+\partial_x^2S^3-S^3\partial_x^2S^+
+\alpha_6(S^+)^2
\right),
\label{eq:class6-eom-plus}\\
\partial_tS^-
&=2\ii\left(
S^3\partial_x^2S^--S^-\partial_x^2S^3
+2\alpha_6(S^3)^2-\alpha_6S^+S^-
\right),
\label{eq:class6-eom-minus}\\
\partial_tS^3
&=\ii\left(
S^-\partial_x^2S^+-S^+\partial_x^2S^-
-2\alpha_6S^+S^3
\right).
\label{eq:class6-eom-three}
\end{align}

The relation to the quantum time evolution can also be checked at
finite lattice spacing.  Let $\sigma_n^i$ denote the Pauli matrix
$\sigma^i$ acting at site $n$, with identity operators on the other
sites, and write
$\boldsymbol\sigma_n=(\sigma_n^1,\sigma_n^2,\sigma_n^3)^{\mathsf T}$.
Taking the lower symbol of the Heisenberg derivative componentwise on
a product of coherent states at a fixed lattice time gives
\begin{align}
\frac{1}{\epsilon^2}
\left(\partial_{t_{\rm lat}}\boldsymbol\sigma_n\right)^\downarrow
={}&-2\boldsymbol S_n\times\Biggl[
\frac{\boldsymbol S_{n-1}+\boldsymbol S_{n+1}-2\boldsymbol S_n}
{\epsilon^2}
\nonumber\\
&\hspace{6em}
-\frac{\alpha_6}{2}N_6
\left(\boldsymbol S_{n-1}+\boldsymbol S_{n+1}\right)
\Biggr].
\label{eq:class6-finite-lattice-spin-flow}
\end{align}
This identity follows from the two bonds adjacent to site $n$ and the
Pauli commutation relations.  Its smooth-field limit is
Eq.~\eqref{eq:class6-vector-eom}, with the time convention
$t=\epsilon^2t_{\rm lat}$.  Applying the same commutator to the
normalized spatial operator gives
\begin{equation}
\left(\partial_{t_{\rm lat}}\widehat L_{6;a,n}\right)^\downarrow
=\epsilon^3\,\partial_tU_6(x_n,t;\lambda)
+\mathcal O(\epsilon^4),
\label{eq:class6-heisenberg-to-Ut}
\end{equation}
where $\partial_t\boldsymbol S$ is given by
Eq.~\eqref{eq:class6-vector-eom}.
Together with the correction identity
\eqref{eq:class6-DeltaV-absorbs-C}, this supplies both sides of the
continuum zero-curvature equation.

To establish equivalence with the equations of motion, we now treat
the time derivatives as independent variables.  Define the
equation-of-motion residual vector
\begin{equation}
\boldsymbol E^{(6)}
:=\partial_t\boldsymbol S
+2\boldsymbol S\times
\left(\partial_x^2\boldsymbol S-\alpha_6N_6\boldsymbol S\right).
\label{eq:class6-residual-vector}
\end{equation}
Write $E^{(6)}_i$ for its Cartesian components and
$E^{(6)}_\pm:=E^{(6)}_1\pm\ii E^{(6)}_2$ for its transverse components.
The curvature of the matrices in
Eqs.~\eqref{eq:class6-spatial-Lax} and
\eqref{eq:class6-time-Lax-compact} is
\begin{equation}
F^{(6)}_{tx}:=\partial_tU_6-\partial_xV_6+[U_6,V_6].
\label{eq:class6-curvature-definition}
\end{equation}
Using only
$\boldsymbol S^{\mathsf T}\boldsymbol S=1$,
$\boldsymbol S^{\mathsf T}\partial_x\boldsymbol S=0$, and
$\boldsymbol S^{\mathsf T}\partial_x^2\boldsymbol S
=-\partial_x\boldsymbol S^{\mathsf T}\partial_x\boldsymbol S$,
we obtain the off-shell identity
\begin{equation}
F^{(6)}_{tx}
=\frac{1}{4\lambda}
\begin{pmatrix}
E^{(6)}_3+2\alpha_6\lambda^2E^{(6)}_+
&E^{(6)}_-+4\alpha_6\lambda^2E^{(6)}_3
-4\alpha_6^2\lambda^4E^{(6)}_+\\
E^{(6)}_+
&-E^{(6)}_3-2\alpha_6\lambda^2E^{(6)}_+
\end{pmatrix}.
\label{eq:class6-offshell-factorization}
\end{equation}
For $\lambda\ne0$, the inverse relations are
\begin{align}
E^{(6)}_+
&=4\lambda\bigl(F^{(6)}_{tx}\bigr)_{21},
\label{eq:class6-residual-plus-from-curvature}\\
E^{(6)}_3
&=4\lambda\left[
\bigl(F^{(6)}_{tx}\bigr)_{11}
-2\alpha_6\lambda^2\bigl(F^{(6)}_{tx}\bigr)_{21}
\right],
\label{eq:class6-residual-three-from-curvature}\\
E^{(6)}_-
&=4\lambda\Bigl[
\bigl(F^{(6)}_{tx}\bigr)_{12}
-4\alpha_6\lambda^2\bigl(F^{(6)}_{tx}\bigr)_{11}
\nonumber\\
&\hspace{4.5em}
+12\alpha_6^2\lambda^4\bigl(F^{(6)}_{tx}\bigr)_{21}
\Bigr].
\label{eq:class6-residual-minus-from-curvature}
\end{align}
Thus $F^{(6)}_{tx}=0$ is equivalent to all three component equations
\eqref{eq:class6-eom-plus}--\eqref{eq:class6-eom-three} for every
$\lambda\ne0$ and every $\alpha_6$.  This equivalence is a local
identity on the unit-spin constraint surface.
